% !PRINTED 540
\weberpage{540}
%<<CONT
Gruppe der Gleichung vorkommt, so ist $C_e$ diese Gruppe. $C_e$ ist aber, wenn $e < m$ ist, intransitiv, und die Gleichung ist reducibel.

Auch auf die cyklischen Gleichungen, deren Grad keine Primzahl ist, lässt sich die oben durchgeführte Reduction anwenden.

Wenn nämlich $m = ef$ irgend eine Zerlegung von $m$ in zwei Factoren ist, so zerfällt die Permutation $\pi^e$ in $e$ Cyklen $\gamma$ von je $f$ Gliedern:
\weq{10}{
\begin{aligned}
\gamma &= (\alpha, \alpha_e, \alpha_{2e} \ldots \alpha_{(f-1)e}), \\
\gamma_1 &= (\alpha_1, \alpha_{e+1}, \alpha_{2e+1} \ldots \alpha_{(f-1)e+1}), \\
&\cdots\cdots\cdots\cdots\cdots\cdots\cdots\cdots \\
% [sic] authority prints c-1 here; mathematical cycle indexing requires e-1.
\gamma_{c-1} &= (\alpha_{e-1}, \alpha_{2e-1}, \alpha_{3e-1} \ldots \alpha_{m-1}).
\end{aligned}
}
Nun bestimmen wir eine zu dem ersten dieser Cyklen gehörige Function $\psi$ und setzen
\weq{11}{
\begin{aligned}
\eta &= \psi(\alpha, \alpha_e, \alpha_{2e} \ldots \alpha_{(f-1)e}), \\
\eta_1 &= \psi(\alpha_1, \alpha_{e+1}, \alpha_{2e+1} \ldots \alpha_{(f-1)e+1}), \\
&\cdots\cdots\cdots\cdots\cdots\cdots\cdots\cdots \\
\eta_{e-1} &= \psi(\alpha_{e-1}, \alpha_{2e-1}, \alpha_{3e-1} \ldots \alpha_{m-1}).
\end{aligned}
}
Diese Grössen sind alle von einander verschieden, aber bei einer cyklischen Permutation ihrer Argumente bleiben sie ungeändert. Durch Anwendung der Permutation $\pi$ gehen die Grössen
\[
\eta, \eta_1, \eta_2 \ldots \eta_{e-1}
\]
cyklisch in einander über, und nach der Voraussetzung sind also ihre cyklischen Functionen, und folglich auch ihre symmetrischen Functionen rational. Sie sind also die Wurzeln einer cyklischen Gleichung $e^{\mathrm{ten}}$ Grades, während $F(x)$ in $e$ Factoren $f^{\mathrm{ten}}$ Grades
\[
F(x) = F(x, \eta)\, F(x, \eta_1) \ldots F(x, \eta_{e-1})
\]
zerfällt, deren jeder eine cyklische Gleichung $f^{\mathrm{ten}}$ Grades für die Wurzeln eines der Cyklen $\gamma$ ergiebt.

Denn setzen wir etwa
\[
F_1 = (x - \alpha)(x - \alpha_e) \ldots (x - \alpha_{(f-1)e}),
\]
so gestattet diese Function die Permutationen der Periode von $\gamma$. Also gestattet sie auch die Permutationen der Gruppe, die nach Adjunction von $\eta$ zur \gesperrt{Galois}'schen Gruppe unserer Gleichung wird, die ja nur aus Potenzen von $\gamma, \gamma_1 \ldots \gamma_{e-1}$ bestehen kann. Es ist daher nach dem Satze von \gesperrt{Lagrange} [\S.\,155, (2)] $F_1$ rational durch $\eta$ darstellbar, also $F_1 = F(x, \eta)$. Die Gleichung
%CONT>>
